\maketitle
\thispagestyle{firstpage}

\begin{abstract}
Physics has temperatures, economics has prices, statistics has regularization strengths, and models of attention have precision weights. Each is a single number that prices a limited resource. This paper asks why such numbers appear whenever a stable descriptive layer forms, and answers within the Six Birds closure calculus. We define a currency relative to a layer by the job it does: bounding what is feasible, adding up along protocols, driving a monotone potential, or pricing a constraint. We then state the currency--constraint principle as a conditional result. If maintaining a higher-level object requires keeping a lower-level ledger within a budget, that ledger becomes a feasibility constraint for the higher layer; if the higher layer then closes by maximizing entropy under the budget, the Gibbs multiplier of the constraint is a nonnegative shadow price. We prove this finite duality, including exactly when the price is positive, zero, or infinite. We also formalize in Lean~4 the single-row duality and the data-processing inequality that prevents coarse observation from creating path-reversal asymmetry. A finite Markov laboratory illustrates each step. Coarse observation never inflates a regularized path audit. The fitted price follows the predicted curve and vanishes once the budget is slack. A resolution ladder exposes more cycle directions (0, 5, 11) than are actually driven (0, 0, 6). A ledger-derived cost predicts held-out transitions better than a scale-matched proxy. Stronger maintenance pricing shrinks packaging defects under an explicit bound. Persistence alone does not create a budget; the principle applies once a resource law is given, and prices then follow whenever the layer closes by maximum entropy under a feasible budget.
\end{abstract}
